\section{Methods}

\sysname represents a design as a directed acyclic graph of contracts, where each node is a component with a contract over its ports. At an internal node, an agent proposes one contract for each subcomponent. The proposal is accepted only if the component contracts are consistent, bounded checks on the actual components find no counterexample, and their composition refines the parent contract. Each subcomponent is then discharged by a back end or, if this does not terminate within a budget, further decomposed. When a node cannot satisfy its contract, it returns a witness to its parent, which revises the decomposition while certified subtrees are preserved. The leaves depend on the domain: a hardware module whose contract a model checker proves directly, a task that one LLM agent completes within its context or a piece of deterministic code, or a continuous-time robot controller certified, for instance, by a control barrier function or a Hamilton-Jacobi reachability analysis. The research is organized in three thrusts.

\thrust{Thrust 1: Decomposition of temporal contracts.} We will study the decomposition of a system contract in the GR(1) fragment~\cite{bloem2012gr1} into realizable component contracts over a given architecture. First, we will identify classes of architectures and contracts for which the existence of a decomposition is decidable, starting from pipeline architectures, for which distributed synthesis is decidable~\cite{pnueli1990distributed}, and use the agent as a search heuristic outside these classes. Second, we will develop sound GR(1) approximations of the quotient. An approximation that strengthens the guarantees of the exact quotient can be solved directly by a synthesis tool, and any controller synthesized for it is correct for the composition. When the approximation is unrealizable, its unrealizable core indicates which terms to weaken in the next approximation. Third, we will handle circular dependencies by restricting the guarantees that can be used as assumptions to safety properties, adopting a circular \gls{ag} rule with an inductive side condition~\cite{mcmillan1999circular,namjoshi2000completeness}, and characterizing the class of decompositions this rule admits. Fourth, as learned clauses prune the search of a SAT solver, we will generalize the unrealizable core of a failed decomposition, via subgraph isomorphism~\cite{xiao2024contrarc}, into a constraint that excludes every candidate assigning an isomorphic component the same terms.

\thrust{Thrust 2: Temporal obligations in long-horizon agents.} Tasks with long horizons, hundreds of tools, and long policy documents require an agent to compact its context and to delegate subtasks, and both steps can lose the rules the agent must follow~\cite{cemri2025mast}; whether one agent or several perform better depends on the structure of the task~\cite{kim2025scaling}. We will keep these rules in a library of contracts in \gls{ltlf}~\cite{degiacomo2013ltlf} stored outside the context window, so that compaction cannot remove them. The library is populated from the instructions and from the context at each compaction, and is extended with contracts mined from execution traces~\cite{lemieux2015texada}, while consistency, vacuity~\cite{maoz2020vacuity}, and well-separation~\cite{maoz2016wellsep} checks reject faulty formalizations. Delegation then becomes a decomposition step: a parent agent passes to each sub-agent a task description, a context, and a set of contracts, either inherited or generated for it, and the engine checks that the composition of the sub-agent contracts refines the parent contract. A central question is how to generate sub-agent contracts that refine the parent by construction, since a sub-agent may need weaker assumptions, stronger guarantees, and conditions on its siblings, such as sending a confirmation only after a payment; the quotient of the parent contract by the contracts of the siblings gives the weakest admissible contract. The contracts also bound the tools exposed to the agent at each step. We will cast the choice between one agent and several as a contract-based architecture exploration problem~\cite{xiao2024contrarc,xiao2026smcp}. Enforcement at run time reduces to monitoring~\cite{xiao2026symbolic} and is outside the scope of this project.

\thrust{Thrust 3: Formal feedback, repair, and reuse.} The contract engine separates a contract algebra independent of the formalism, contract-level checks reduced to validity or satisfiability, and analyses specific to a formalism, such as unrealizable cores for GR(1). Every failed check returns a counterexample trace or a minimal core identifying the terms responsible for the failure. A \gls{mcp} server and an agent skill expose all the operations to LLM agents, while no LLM is part of the engine. A property is claimed only after the back-end result has been read and the vacuity checks have passed, and bounded or timed-out runs count as undetermined. Refinement failures are first repaired mechanically using the quotient, and the agent is invoked only when this repair does not succeed. A certified component can be reused in another design after its assumptions are rechecked. Finally, we will measure the coverage of certified contracts by injecting faults into the components.

\thrust{Applications and preliminary results.} We will evaluate \sysname in three domains. In hardware verification, component contracts are \glspl{sva} with ghost-state models, verified with Cadence Jasper and, in this project, the open-source SymbiYosys. Our prototype certified an end-to-end in-order delivery property of an asynchronous link from the HierSVA benchmark with 25 Jasper runs and about 16,000 LLM input tokens, which flat Jasper did not decide in the same time. The decomposition found at FIFO depth 8 certified depths 16 to 512 with no further LLM calls. An agent with the same tools and no protocol used 24 times as many tokens without completing the proof. In multi-robot control, ACIES~\cite{xiao2026acies} decomposes mission requirements into GR(1) contracts and synthesizes controllers; in a three-level decomposition with 42 nodes, the contract library closed 26 of the 39 leaves. For LLM agents, we will start from a single-agent contract-based harness and add context compaction, contract extraction, trace mining, and contract-checked delegation one at a time, on long-horizon tasks from AppWorld~\cite{trivedi2024appworld} and policy-heavy benchmarks, against the best summarization-based compaction as baseline.
